\section{Building The Corpora}

The natural first step was to look at other corpora that aggregate similar data.  

\textbf{Europarl} dataset \cite{koehn_europarl_2005} aggregates proceedings of the European Parliament in twenty one different European languages\footnote{All 21 Europarl languages: Romanic (French, Italian, Spanish, Portuguese, Romanian), Germanic (English, Dutch, German, Danish, Swedish), Slavik (Bulgarian, Czech, Polish, Slovak, Slovene), Finni-Ugric (Finnish, Hungarian, Estonian), Baltic (Latvian, Lithuanian), and Greek.} with the goal of generating sentence aligned text for statistical machine learning translation systems. The overall structure is akin to a collection of plaintext with minimal tagging.

\textbf{ItaParlCorpus} \cite{cova_new_2025}, on the other hand, is a comprehensive, annotated, and machine-readable database of Italy's parliamentary speeches spanning from 1948 to 2022. Each speech is identified by a precise date, the role of the speaker and its party name and family, and by the legislature number. It is not, whoever, very rich in terms of linguistic processing tagging.

\textbf{IMPAQTS} \cite{cominetti_impaqts_2024} is a multimodal corpus of around 2.65 million tokens including 1,500 speeches uttered by 150 prominent politicians spanning from 1946 to 2023. For each speaker, the corpus contains 4 parliamentary speeches, 2 rallies, 1 party assembly, and 3 statements (in person or broadcasted). The goal is to annotate political rhetoric that hides a deeper meaning, for example by employing sarcasm. There are four main types of annotation: implication (conventional, generalized, etc), presupposition (e.g., pragmatic), vagueness (syntactic, semantic, metaphoric), and topicalization (either syntactic or prosodic).

\textbf{Parola di Leader}, and specifically the corpus LP4 \cite{giuliano_il_2019}, is a corpus of more than five million occurrences taken from the stenographic archives of the Chamber of Deputies of Italy\footnote{Chamber of Deputies official online public archive: \url{https://legislatureprecedenti.camera.it/}.}, between 1948 (first legislature of the Italian Republic) and 2011 (excluding Monti Government, the $61^{st}$ legislature).

They selected a list of thirty "highly influential politicians"\footnote{The full list of politicians follows: Giorgio Almirante, Giuliano Amato, Giulio Andreotti, Enrico Berlinguer, Silvio Berlusconi, Pier Luigi Bersani, Fausto Bertinotti, Rosy Bindi, Emma Bonino, Umberto Bossi, Pier Ferdinando Casini, Francesco Cossiga, Bettino Craxi, Massimo D'Alema, Alcide De Gasperi, Ciriaco De Mita, Antonio Di Pietro, Amintore Fanfani, Gianfranco Fini, Ugo La Malfa, Aldo Moro, Pietro Nenni, Achille Occhetto, Marco Pannella, Romano Prodi, Giuseppe Saragat, Giovanni Spadolini, Palmiro Togliatti, Valter Veltroni, Nichi Vendola.} and then for each of them they selected some speeches. Overall, they collected 1877 distinct speeches.

The corpus could be considered a corpora made of thirty distinct corpus, one for each politician. Each of these has, among others, the following features: name of the politician, role of politician (for example Prime Minister), whether the politician is the major party or in the opposition, number of legislature and number of seating of the day, and date of the seating, argument of the speech, name of the current Government (e.g., \textit{Berlusconi I}), and lastly a reference to the file that holds the actual speech.

There is also a corpus for the vocabulary that holds, for each word, the total number of occurrences, a Part-Of-Speech\footnote{In grammar, a part of speech or part-of-speech (abbreviated as POS or PoS) is a category of words (or, more generally, of lexical items) that have similar grammatical properties \cite{payne_describing_1997}.} tag, its lemma, a count for the number of occurrences for each politicians, major party and opposition, and lastly the imprinting tag (e.g., masculine singular).

Lastly, \textbf{ParlaMint} \cite{erjavec_parlamint_2023,erjavec_parlamint_2025} is a compiled comparable parliamentary corpora for a number of countries and languages. ParlaMint corpora are interoperable, i.e. encoded to a very constrained common ParlaMint schema, a specialization of the Parla-CLARIN\footnote{Prla-CLARIN, a TEI schema for corpora of parliamentary proceedings: \url{https://clarin-eric.github.io/parla-clarin/}.} recommendations, which are a customization of the TEI Guidelines\footnote{P5 TEI Guidelines: \url{https://tei-c.org/guidelines/p5/}.}.

It is a collection of parliamentary seating split into utterances. An example of two utterances could be \textit{Politician: "I ask permission to speak."}, \textit{President of the Senate: "Permission granted."}. For each utterance there are, among others, the following metadata: ID for both the whole seating and the utterance itself, a title of the seating, a date of the seating, the State Body that where it did take place (e.g., the Senate of the Republic), the name of the speaker and its affiliation, as well as the name of its coalition. 

Speeches are collected before splitting in their own dataset, and for each speech we have, among others, a unique ID and a date, the type of speech (e.g., \textit{Interview} or \textit{Rally}), the venue channel (e.g., \textit{BBC News}), the unique ID of the speaker and its role (e.g., \textit{host} or \textit{main speaker}), the text of the speech itself and its type (e.g., \textit{Question}, \textit{Answer}, \textit{Interruption}) and whether there have been some \textit{Incidents} and their type.

Incidents are breaks in the flow of the conversation, and they can be of different types\footnote{Full list of incident types: "action", "incident", "leaving", "entering", "break", "pause", "sound", "editorial".}, for example \textit{"leaving"} or \textit{"break"}. If an incident is of the kinesic type, it can be further subdivided ito different kinesic types, such as \textit{"applause"} or \textit{"laughter"}\footnote{Full list of kinesic types: "kinesic", "applause", "ringing", "signal", "playback", "gesture", "smiling", "laughter", "snapping", "noise".}. If the incident is of the vocal type, it can be further subdivided ito different vocal types, such as \textit{"question"} or \textit{"shouting"}\footnote{Full list of vocal types: "greeting", "question", "clarification", "speaking", "interruption", "exclamat", "laughter", "shouting", "murmuring", "noise", "signal".}.

Speakers and Parties have also their own datasets, with unique IDs, names, affiliations, roles and the like.

Overall the taxonomy of ParlaMint is strongly codified, which make the interpretation of the dataset easier and machine-readable. All data is saved also in XML format.

\subsection{My Own Corpora}

Ultimately, I did not manage to find any corpus that I can directly integrate or use for my purposes, but I took inspiration from ParlaMint for the overall structure. 
Of course at the time of writing I still have not started compiling it and therefore the structure may be subjected of minor changes.